% Lineare Funktionen · Anwendung · Prüfungs-Fokus MSA – bank/lineare-funktionen/e5.jsonl (zusammenbau v0.6, Prüfungs-Fokus)
\einheitenkopf[e5]{Einheit 5 · Anwendungen}
\zweigzeile{ab Kl. 8 · P10 ×6}
\setcounter{aufgabe}{0}

% Anlauf (ohne Prüfkennung)
\begin{aufgabe}{Anwendung – Anlauf}
\begin{teile}
\teil Markiere: Unterstreiche den Anfangswert, kreise die Änderung je Einheit ein. Ein Handwerker berechnet 45 € für die Anfahrt und 38 € je Arbeitsstunde.
\teil Welche Gleichung passt? Ein Fahrradverleih nimmt 8 € Grundgebühr und 3 € je Stunde. y ist der Preis in €, x die Anzahl der Einheiten. Kreuze an. \\ \kreuz{$y = 8x + 3$} \\ \kreuz{$y = 3x + 8$} \\ \kreuz{$y = 11x$} \\ \kreuz{$y = 3x - 8$}
\teil Stelle eine Gleichung auf: Ein Kino-Abo kostet einmalig 10 € und 6 € je Film. y sind die Kosten in €, x die Anzahl der Filme. y = \leerfeld
\end{teile}
\end{aufgabe}

\begin{aufgabe}{Anwendung – Prüfungsaufgaben}
\begin{teile}
\steil Welches Mietauto? Angebot 1: 32 € Miete pro Tag, 300 km insgesamt frei, jeder weitere km 0{,}30 €. Angebot 2: einmalig 95 € für eine Woche, jeder km 0{,}12 €. Jana mietet 4 Tage und fährt 420 km. Berechne die Gesamtkosten beider Angebote, entscheide, welches günstiger ist, und stelle für Angebot 2 eine Gleichung für die Kosten y (in €) bei x km für eine Woche auf. \hfill (P23*) \\ Angebot 1: \leerfeld[€] \quad Angebot 2: \leerfeld[€] \quad y = \leerfeld
\steil Welches Mietauto? Angebot 1: 29 € Miete pro Tag, 250 km insgesamt frei, jeder weitere km 0{,}40 €. Angebot 2: einmalig 110 € für eine Woche, jeder km 0{,}10 €. Jana mietet 3 Tage und fährt 380 km. Berechne die Gesamtkosten beider Angebote, entscheide, welches günstiger ist, und stelle für Angebot 2 eine Gleichung für die Kosten y (in €) bei x km für eine Woche auf. \hfill (P23*) \\ Angebot 1: \leerfeld[€] \quad Angebot 2: \leerfeld[€] \quad y = \leerfeld
\teil Welches Busunternehmen? Nord verlangt 150 € Grundpreis und 1{,}20 € je km, Süd 1{,}80 € je km ohne Grundpreis. Das Ziel ist 120 km entfernt, der Bus fährt hin und zurück. Welches Unternehmen ist günstiger? Bleibt es dabei, wenn am Ziel noch ein Ausflug von 80 km dazukommt? Begründe. \hfill (P16)
\teil Welches Busunternehmen? Adler verlangt 90 € Grundpreis und 0{,}80 € je km, Biber 1{,}10 € je km ohne Grundpreis. Das Ziel ist 130 km entfernt, der Bus fährt hin und zurück. Welches Unternehmen ist günstiger? Bleibt es dabei, wenn am Ziel noch ein Ausflug von 60 km dazukommt? Begründe. \hfill (P16)
\teil Welche Gleichung gehört zu welchem Sachverhalt? Sachverhalt 1: E-Scooter: 1 € fürs Entsperren und 25 Cent je Minute. Sachverhalt 2: Schwimmbad: 25 € Jahresgebühr und 1 € je Besuch. y ist der Preis in €, x die Anzahl der Einheiten. Ordne jedem Sachverhalt eine der Gleichungen $y = 25x + 1$, $y = x + 25$, $y = 0{,}25x + 1$ zu. \hfill (P21) \\ Sachverhalt 1: \leerfeld \quad Sachverhalt 2: \leerfeld
\teil Welche Gleichung gehört zu welchem Sachverhalt? Sachverhalt 1: Parkhaus: 2 € bei der Einfahrt und 60 Cent je Stunde. Sachverhalt 2: Sportverein: 60 € Aufnahmegebühr und 2 € je Training. y ist der Preis in €, x die Anzahl der Einheiten. Ordne jedem Sachverhalt eine der Gleichungen $y = 60x + 2$, $y = 2x + 60$, $y = 0{,}60x + 2$ zu. \hfill (P21) \\ Sachverhalt 1: \leerfeld \quad Sachverhalt 2: \leerfeld
\steil Wann ist das Ziel erreicht? Lina hat 640 € auf dem Konto und zahlt jeden Monat 45 € ein. y ist das Guthaben in €, x die Anzahl der Monate. Stelle eine Gleichung für das Guthaben auf und berechne, nach wie vielen Monaten Lina mindestens 1\,500 € hat. \hfill (P22*) \\ y = \leerfeld \quad nach \leerfeld Monaten
\steil Wann ist das Ziel erreicht? Jan hat 1\,150 € auf dem Konto und zahlt jeden Monat 70 € ein. y ist das Guthaben in €, x die Anzahl der Monate. Stelle eine Gleichung für das Guthaben auf und berechne, nach wie vielen Monaten Jan mindestens 2\,000 € hat. \hfill (P22*) \\ y = \leerfeld \quad nach \leerfeld Monaten
\steil Funktionsgleichung: Ein Umzugsunternehmen berechnet 180 € Grundpreis und 1{,}20 € je km. Gib die Funktionsgleichung für den Preis y (in €) bei x km an. \hfill (P16*) \\ y = \leerfeld
\steil Funktionsgleichung: Ein Reisebus kostet 250 € Grundpreis und 2{,}40 € je km. Gib die Funktionsgleichung für den Preis y (in €) bei x km an. \hfill (P16*) \\ y = \leerfeld
\teil Kontostand nach einem Jahr? Ole hat 830 € auf dem Konto und zahlt jeden Monat 35 € ein. Wie viel Geld hat er nach einem Jahr? \hfill (P22) \\ \leerfeld[€]
\teil Bäume nach 6 Jahren? Eine Stadt hat 2\,340 Straßenbäume und pflanzt jedes Jahr 85 neue dazu. Wie viele Straßenbäume hat sie nach 6 Jahren, wenn keiner gefällt wird? \hfill (P22) \\ \leerfeld Bäume
\teil Werte ergänzen: Ein Wassertank wird gleichmäßig geleert. Die Tabelle zeigt die Wassermenge. Trage die Wassermenge zu Beginn und nach 40 Minuten ein. \hfill (P21)

\wertetabelle[,270,240,210,180,]{Zeit in min}{Wasser in l}{0,5,10,15,20,40}
\teil Werte ergänzen: Ein Akku wird gleichmäßig geladen. Die Tabelle zeigt den Ladestand. Trage den Ladestand zu Beginn und nach 90 Minuten ein. \hfill (P21)

\wertetabelle[,26,32,38,44,]{Zeit in min}{Ladung in \%}{0,10,20,30,40,90}
\teil Welcher Graph gehört zu welcher Firma? Firma Blitz verlangt 40 € Grundpreis und 1{,}20 € je km, Firma Flott 1{,}60 € je km ohne Grundpreis. Ordne die Graphen I und II den Firmen zu und begründe deine Wahl für Blitz. \hfill (P16) \\

\begin{ksys}[xmin=0,xmax=150,ymin=0,ymax=250,xstep=25,ystep=50,xlabel=x in km,ylabel=y in €,ablesen] \gerade{1.6}{0}{I} \gerade{1.2}{40}{II} \end{ksys}
Blitz: \leerfeld \quad Flott: \leerfeld
\teil Welcher Graph gehört zu welchem Verein? Verein Hoch verlangt 30 € Aufnahmegebühr und 5 € je Monat, Verein Weit 8 € je Monat ohne Aufnahmegebühr. Ordne die Graphen I und II den Vereinen zu und begründe deine Wahl für Hoch. \hfill (P16) \\

\begin{ksys}[xmin=0,xmax=12,ymin=0,ymax=100,xstep=1,ystep=10,xlabel=x in Monaten,ylabel=y in €,ablesen] \gerade{5}{30}{I} \gerade{8}{0}{II} \end{ksys}
Hoch: \leerfeld \quad Weit: \leerfeld
\teil Welches Diagramm passt? Angebot 1: 32 € Miete für einen Tag, 300 km frei, jeder weitere km 0{,}30 €. Angebot 2: einmalig 95 €, jeder km 0{,}12 €. Die Graphen A bis D zeigen Kosten y (in €) für einen Tag bei x km. Gib für jedes Angebot den passenden Graphen an. \hfill (P23) \\

\begin{ksys}[xmin=0,xmax=600,ymin=0,ymax=200,xstep=100,ystep=20,xlabel=x in km,ylabel=y in €,ablesen] \funktion{max(32,32+0.3*(\x-300))}{A} \gerade{0.12}{0}{B} \gerade{0.12}{95}{C} \gerade{0.3}{32}{D} \end{ksys}
Angebot 1: \leerfeld \quad Angebot 2: \leerfeld
\teil Welches Diagramm passt? Angebot 1: 25 € Miete für einen Tag, 200 km frei, jeder weitere km 0{,}40 €. Angebot 2: einmalig 80 €, jeder km 0{,}15 €. Die Graphen A bis D zeigen Kosten y (in €) für einen Tag bei x km. Gib für jedes Angebot den passenden Graphen an. \hfill (P23) \\

\begin{ksys}[xmin=0,xmax=500,ymin=0,ymax=160,xstep=100,ystep=20,xlabel=x in km,ylabel=y in €,ablesen] \gerade{0.4}{25}{A} \gerade{0.15}{80}{B} \gerade{0.15}{0}{C} \funktion{max(25,25+0.4*(\x-200))}{D} \end{ksys}
Angebot 1: \leerfeld \quad Angebot 2: \leerfeld
\end{teile}
\end{aufgabe}

\medskip
\uebersichtskasten{Sachverhalt → Gleichung: n = Anfangswert (Grundgebühr), m = Änderung je Einheit (Preis je km). \\ 3,50 € Grundgebühr, 2 € je km: K(x) = 2x + 3,5}
